% ============================================================
%  Informe tecnico — EL-5409 Laboratorio de Control Automatico
%  Laboratorio 1: Model-Based Design para la planta termica TCLab
%  Formato IEEE, dos columnas. Compilar con pdfLaTeX + BibTeX.
% ============================================================
\documentclass[conference,a4paper,10pt]{IEEEtran}

% ============================================================
%  Preambulo — Informe IEEE (dos columnas) · Lab 1 EL-5409
% ============================================================

\usepackage[T1]{fontenc}
\usepackage[utf8]{inputenc}
\usepackage[spanish,es-tabla,es-noquoting]{babel}
\usepackage{amsmath,amssymb}
\usepackage{graphicx}
\usepackage{booktabs}
\usepackage{array}
\usepackage{enumitem}
\usepackage{listings}
\usepackage{xcolor}
\usepackage{balance}
\usepackage{xurl}          % permite cortar URLs largas en la bibliografia
\usepackage[hidelinks,unicode]{hyperref}

% IEEEtran define \IEEEproof etc.; evitamos choques de babel
\addto\captionsspanish{\renewcommand{\tablename}{TABLA}}

% ---------- Unidades ----------
\usepackage{siunitx}
\sisetup{
  output-decimal-marker = {,},
  group-separator       = {\,},
  per-mode              = symbol,
  detect-all            = true
}

% ---------- Graficas ----------
\usepackage{pgfplots}
\usepackage{pgfplotstable}
\pgfplotsset{compat=1.18}
\usepgfplotslibrary{groupplots}
\usetikzlibrary{matrix}

% Ruta a la carpeta de datos
\newcommand{\G}{../Graficas}

\definecolor{cAzul}   {HTML}{1F4E79}
\definecolor{cNaranja}{HTML}{C55A11}
\definecolor{cVerde}  {HTML}{2E7D32}
\definecolor{cGris}   {HTML}{6E6E6E}
\definecolor{cRojo}   {HTML}{A61C1C}

% Estilo base: pensado para una columna IEEE
\pgfplotsset{
  estilo/.style={
    width            = \columnwidth,
    height           = 4.1cm,
    scale only axis  = false,
    grid             = major,
    grid style       = {gray!20, line width=0.25pt},
    tick label style = {font=\scriptsize},
    label style      = {font=\scriptsize},
    legend style     = {font=\scriptsize, draw=gray!45, fill=white,
                        fill opacity=0.85, text opacity=1,
                        inner sep=1.5pt, row sep=-1pt},
    legend cell align = left,
    every axis plot/.append style = {line width=0.7pt},
    scaled ticks = false,
    tick align   = outside,
    max space between ticks = 40pt,
  },
  ancho/.style = {width=0.47\textwidth, height=4.3cm},
  aligera/.style = {each nth point=#1, filter discard warning=false,
                    unbounded coords=discard},
  aligera/.default = 4,
}

% ---------- Codigo fuente ----------
\lstdefinestyle{estiloC}{
  language=C, basicstyle=\ttfamily\scriptsize,
  keywordstyle=\color{cAzul}\bfseries,
  commentstyle=\color{cGris}\itshape,
  frame=single, rulecolor=\color{gray!40},
  breaklines=true, showstringspaces=false, columns=fullflexible,
  xleftmargin=0.6em, aboveskip=4pt, belowskip=2pt
}
\lstset{style=estiloC}

% ---------- Marcas de trabajo pendiente ----------
% Cambie a \mostrarpendientesfalse antes de entregar
\newif\ifmostrarpendientes
\mostrarpendientestrue
\newcommand{\pendiente}[1]{%
  \ifmostrarpendientes
    \par\noindent\fcolorbox{cRojo}{red!4}{%
      \parbox{0.95\columnwidth}{\scriptsize\color{cRojo}\textbf{PENDIENTE:} #1}}\par\smallskip
  \fi}

% ---------- Atajos ----------
\newcommand{\Tref}{T_{\mathrm{ref}}}
\newcommand{\Tamb}{T_{\mathrm{amb}}}
\newcommand{\Qcmd}{Q_{\mathrm{cmd}}}
\newcommand{\Gpnom}{G_{p,\mathrm{nom}}}
\newcommand{\Kp}{K_p}
\newcommand{\Ki}{K_i}
\newcommand{\Ti}{T_i}


\begin{document}

\title{Diseño basado en modelos de un controlador PI\\
de temperatura para la planta térmica TCLab}

\author{%
\IEEEauthorblockN{Abner López Méndez, 2022273075 \qquad Jordi Segura Chinchilla, 2022240646}
\IEEEauthorblockA{Escuela de Ingeniería Electrónica, Tecnológico de Costa Rica\\
EL-5409 --- Laboratorio de Control Automático, II Semestre 2026\\
Profesor: Ing. Juan Luis Guerrero Fernández, PhD}}

\maketitle

\begin{abstract}
Se documenta el desarrollo de un sistema de control de temperatura para la planta térmica TCLab
siguiendo un flujo de diseño basado en modelos. A partir de un balance de energía se obtuvo un
modelo no lineal y su linealización alrededor de $Q_1 = \SI{50}{\percent}$. Dos experimentos sobre
la planta física, un escalón y una secuencia multinivel, permitieron identificar sendos modelos de
primer orden con tiempo muerto; la validación cruzada seleccionó como nominal el modelo obtenido
del ensayo escalón, $\Gpnom(s) = 0{,}979\,e^{-8{,}00 s}/(179{,}1 s + 1)$, con un ajuste del
\SI{88.7}{\percent} sobre datos no empleados en su identificación. Las reglas SIMC con
$\tau_c = 2\theta$ produjeron un controlador PI con $\Kp = \SI{7.62}{\percent\per\celsius}$ y
$\Ti = \SI{96}{\second}$, que en simulación Model-in-the-Loop alcanza un tiempo de establecimiento
de \SI{205}{\second} sin sobreimpulso al incorporar anti-windup por \emph{clamping}, frente a
\SI{392}{\second} y \SI{30.1}{\percent} de sobreimpulso sin él. El código generado con Embedded
Coder reprodujo exactamente la respuesta del modelo tanto en Software-in-the-Loop como en
Processor-in-the-Loop sobre un Arduino Uno R4 Minima. La implementación final sobre la TCLab,
ejecutada en el Arduino Leonardo mediante Monitor \& Tune, alcanzó la referencia de
\SI{60}{\celsius} con un error estacionario de \SI{0.04}{\celsius} y un tiempo de establecimiento
de \SI{298}{\second}, con una ganancia estática aparente cerca de un \SI{15}{\percent} menor que la
del modelo nominal.
\end{abstract}

\begin{IEEEkeywords}
TCLab, diseño basado en modelos, identificación FOPDT, SIMC, anti-windup, SIL, PIL, Embedded Coder.
\end{IEEEkeywords}

\section{Introducción}

La Temperature Control Lab (TCLab) es una planta térmica didáctica montada sobre un Arduino
Leonardo. Dispone de dos transistores que actúan como calentadores y de dos sensores de temperatura
adheridos a ellos; la potencia de cada calentador se comanda como un porcentaje de su valor máximo
y la temperatura se mide en grados Celsius. En este trabajo se utilizó únicamente el primer par
calentador--sensor, de modo que la variable manipulada es $Q_1$, expresada en porcentaje, y la
variable controlada es $T_1$.

El problema de control consiste en llevar $T_1$ desde la temperatura ambiente hasta una referencia
de \SI{60}{\celsius} y mantenerla en ese valor, respetando el límite de potencia
$Q_1 \in [0, 90]\,\si{\percent}$ impuesto para proteger el equipo. Los requisitos de desempeño
establecidos fueron un tiempo de establecimiento $t_s \leq \SI{600}{\second}$ y un sobreimpulso
$M_p \leq \SI{5}{\percent}$.

El objetivo del trabajo no se limita a obtener un controlador que funcione en simulación, sino a
recorrer un flujo completo de diseño basado en modelos (\emph{Model-Based Design}, MBD). En este
enfoque el modelo ejecutable es el eje del desarrollo, y entre el diseño matemático y el sistema
físico se intercalan etapas sucesivas de verificación que comprueban aspectos distintos de la
solución \cite{mw_hil}. El flujo seguido fue

\begin{center}
modelado $\rightarrow$ identificación $\rightarrow$ diseño $\rightarrow$ MIL $\rightarrow$
código $\rightarrow$ SIL $\rightarrow$ PIL $\rightarrow$ implementación física.
\end{center}

Cada etapa acota el riesgo de la siguiente. La evaluación \emph{Model-in-the-Loop} (MIL) comprueba
el desempeño del algoritmo contra un modelo de la planta; \emph{Software-in-the-Loop} (SIL)
verifica que el código generado se comporte como el modelo del que proviene;
\emph{Processor-in-the-Loop} (PIL) verifica que ese código conserve su comportamiento al ejecutarse
sobre el procesador objetivo; y solo entonces el controlador se cierra sobre la planta real. El
documento sigue ese mismo orden.

\section{Modelado e identificación del sistema}

\subsection{Modelo físico}

El conjunto formado por el transistor calentador y el sensor se representó como una única masa a
temperatura uniforme. Bajo esa hipótesis, y considerando pérdidas por convección natural y por
radiación hacia el entorno, el balance de energía conduce a \cite{hedengren_modeling}

\begin{equation}
m c_p \frac{dT}{dt} = \alpha Q - UA\,(T - T_a) - \epsilon\sigma A\,(T^4 - T_\infty^4),
\label{eq:nolineal}
\end{equation}

donde $\alpha$ convierte el porcentaje de potencia comandado en vatios entregados al sistema, $UA$
es el coeficiente global de convección, y $\epsilon$ y $\sigma$ son la emisividad y la constante de
Stefan--Boltzmann. Las hipótesis adoptadas fueron masa concentrada y temperatura uniforme,
convección predominantemente natural, potencia del calentador proporcional al comando, y
$T_a \approx T_\infty$ por suponer el equipo y el recinto inicialmente en equilibrio térmico. Las
temperaturas del término radiativo se manejaron en kelvin.

La no linealidad de \eqref{eq:nolineal} proviene exclusivamente del término $T^4$. Linealizando
alrededor de un punto de operación $(\bar T, \bar Q)$ y definiendo las variables de desviación
$\Delta T = T - \bar T$ y $\Delta Q = Q - \bar Q$, se obtiene \cite{hedengren_linearize}

\begin{equation}
\frac{d\,\Delta T}{dt} = \gamma\,\Delta T + \beta\,\Delta Q,
\label{eq:lineal}
\end{equation}

\vspace{-3mm}
\begin{equation}
\gamma = -\frac{UA + 4\epsilon\sigma A \bar T^{3}}{m c_p},
\qquad
\beta = \frac{\alpha}{m c_p},
\label{eq:gammabeta}
\end{equation}

cuya función de transferencia es de primer orden,

\begin{equation}
G(s) = \frac{K}{\tau s + 1},
\qquad \tau = -\frac{1}{\gamma},
\qquad K = -\frac{\beta}{\gamma}.
\label{eq:g1orden}
\end{equation}

En \eqref{eq:lineal} el coeficiente $\gamma$, dado por \eqref{eq:gammabeta}, reúne las dos vías de
pérdida, y es el término radiativo el que lo hace depender del punto de operación. El punto empleado fue
$\bar Q = \SI{50}{\percent}$; los valores de $\tau$ y $K$ que resultan de \eqref{eq:g1orden} se
resumen en la Tabla~\ref{tab:fisico}.

\begin{table}[!t]
\caption{Modelo físico linealizado en $\bar Q = \SI{50}{\percent}$ y error de la aproximación
lineal frente al modelo no lineal.}
\label{tab:fisico}
\centering
\footnotesize
\begin{tabular}{@{}lr@{}}
\toprule
\textbf{Magnitud} & \textbf{Valor} \\
\midrule
Temperatura de equilibrio $\bar T$        & \SI{60.47}{\celsius} \\
Pérdida por convección en equilibrio      & \SI{0.2188}{\watt} (\SI{43.8}{\percent}) \\
Pérdida por radiación en equilibrio       & \SI{0.2811}{\watt} (\SI{56.2}{\percent}) \\
$\gamma$                                  & \SI{-7.548e-3}{\per\second} \\
$\beta$                                   & \SI{5.000e-3}{\celsius\per\percent\per\second} \\
Constante de tiempo $\tau$                & \SI{132.5}{\second} \\
Ganancia estática $K$                     & \SI{0.6625}{\celsius\per\percent} \\
\midrule
\multicolumn{2}{@{}l}{\textit{Error del modelo lineal frente al no lineal}} \\
$\Delta Q = \SI{5}{\percent}$: error final / RMSE   & \SI{0.029}{\celsius} / \SI{0.022}{\celsius} \\
$\Delta Q = \SI{40}{\percent}$: error final / RMSE  & \SI{1.744}{\celsius} / \SI{1.337}{\celsius} \\
\bottomrule
\end{tabular}
\end{table}

Los dos últimos renglones de la Tabla~\ref{tab:fisico} delimitan el alcance de la aproximación
lineal. Ante una perturbación de \SI{5}{\percent} el error no llega a \SI{0.03}{\celsius}, mientras
que al multiplicar la amplitud por ocho el error crece casi sesenta veces. Este crecimiento no es
proporcional porque el término omitido es de cuarto orden en la temperatura, y en el punto de
operación la radiación ya aporta el \SI{56.2}{\percent} de las pérdidas totales. Es decir, el
modelo lineal es útil para excursiones pequeñas alrededor del punto de operación, pero no para
describir el calentamiento completo desde temperatura ambiente, que es justamente el transitorio
que debe manejar el controlador. Esta limitación motivó obtener el modelo de diseño de forma
experimental en lugar de usar los parámetros físicos.

\subsection{Identificación experimental}

Se realizaron dos experimentos sobre la unidad TCLab n.\textsuperscript{o}~4, con temperatura
ambiente de \SI{24}{\celsius} y periodo de muestreo $T_s = \SI{1}{\second}$. El primero aplicó un
escalón de $Q_1$ de \SIrange{0}{50}{\percent} y se registró durante \SI{1860}{\second}
(\num{1861} muestras), partiendo de $T_1 = \SI{24.78}{\celsius}$. El segundo aplicó una secuencia
multinivel con escalones sucesivos de \SIrange{10}{50}{\percent} durante \SI{3520}{\second}
(\num{3521} muestras), partiendo de $T_1 = \SI{25.27}{\celsius}$. El ensayo escalón concentra la
información en una sola transición amplia, mientras que el multinivel recorre varios niveles de
potencia y permite observar si la ganancia del proceso depende de la región de operación.

Ambos conjuntos se ajustaron a un modelo de primer orden con tiempo muerto (FOPDT)
\cite{hedengren_fopdt},

\begin{equation}
G_p(s) = \frac{K\,e^{-\theta s}}{\tau s + 1},
\label{eq:fopdt}
\end{equation}

expresado en variables de desviación, con $[K] = \si{\celsius\per\percent}$ y
$[\tau] = [\theta] = \si{\second}$. Los parámetros de \eqref{eq:fopdt} se obtuvieron por dos vías: el procedimiento de
mínimos cuadrados sobre la respuesta simulada empleado en la práctica, y la función
\texttt{procest} de \emph{System Identification Toolbox} con estructura \texttt{P1D}. La
Tabla~\ref{tab:fopdt} recoge los cuatro modelos y la Fig.~\ref{fig:experimentos} muestra las
señales experimentales junto con la respuesta de los modelos identificados.

\begin{table}[!t]
\caption{Parámetros FOPDT identificados por ensayo y por método.}
\label{tab:fopdt}
\centering
\footnotesize
\begin{tabular}{@{}llS[table-format=1.3]S[table-format=3.1]S[table-format=1.2]S[table-format=2.1]@{}}
\toprule
\textbf{Ensayo} & \textbf{Método} & {$K$ (\si{\celsius\per\percent})} & {$\tau$ (\si{\second})}
& {$\theta$ (\si{\second})} & {Ajuste (\si{\percent})} \\
\midrule
Escalón    & Práctica              & 0.979 & 179.1 & 8.00 & 93.0 \\
Escalón    & \texttt{procest} P1D  & 0.979 & 178.7 & 8.36 & 93.0 \\
Multinivel & Práctica              & 1.036 & 189.5 & 2.00 & 94.5 \\
Multinivel & \texttt{procest} P1D  & 1.036 & 189.2 & 2.80 & 94.5 \\
\bottomrule
\end{tabular}
\end{table}

\begin{figure*}[!t]
\centering
\begin{tikzpicture}
\begin{groupplot}[group style={group size=2 by 2, horizontal sep=1.3cm, vertical sep=0.9cm},
                  estilo, ancho]
\nextgroupplot[ylabel={$Q_1$ (\si{\percent})}, xticklabels={}, ymin=-5, ymax=60,
               title={\scriptsize Ensayo escalón}, title style={yshift=-2pt}]
  \addplot[cNaranja, aligera=3] table[col sep=comma, x=Tiempo_s, y=Q1_pct] {\G/IdentificacionSE.csv};
\nextgroupplot[ylabel={$Q_1$ (\si{\percent})}, xticklabels={}, ymin=-5, ymax=60,
               title={\scriptsize Ensayo multinivel}, title style={yshift=-2pt}]
  \addplot[cNaranja, aligera=5] table[col sep=comma, x=Tiempo_s, y=Q1_pct] {\G/IdentificacionSM.csv};
\nextgroupplot[ylabel={$T_1$ (\si{\celsius})}, xlabel={Tiempo (\si{\second})}, legend pos=south east]
  \addplot[cAzul, aligera=3] table[col sep=comma, x=Tiempo_s, y=T1_experimental_C] {\G/IdentificacionSE.csv};
  \addplot[cVerde, dashed, aligera=3] table[col sep=comma, x=Tiempo_s, y=T1_modelo_C] {\G/IdentificacionSE.csv};
  \legend{Experimental, FOPDT}
\nextgroupplot[ylabel={$T_1$ (\si{\celsius})}, xlabel={Tiempo (\si{\second})}, legend pos=south east]
  \addplot[cAzul, aligera=5] table[col sep=comma, x=Tiempo_s, y=T1_experimental_C] {\G/IdentificacionSM.csv};
  \addplot[cVerde, dashed, aligera=5] table[col sep=comma, x=Tiempo_s, y=T1_modelo_C] {\G/IdentificacionSM.csv};
  \legend{Experimental, FOPDT}
\end{groupplot}
\end{tikzpicture}
\caption{Señales de los dos experimentos de identificación y respuesta de los modelos FOPDT
ajustados a cada uno. Izquierda: escalón de $Q_1$ a \SI{50}{\percent}. Derecha: secuencia
multinivel de \SIrange{10}{50}{\percent}.}
\label{fig:experimentos}
\end{figure*}

Los dos métodos coinciden en $K$ y $\tau$ hasta la tercera cifra significativa, lo que respalda la
identificación. La única diferencia apreciable está en $\theta$: el procedimiento de la práctica
cuantiza el tiempo muerto a un número entero de periodos de muestreo, mientras que
\texttt{procest} lo estima como variable continua. Las diferencias resultantes, \SI{0.36}{\second}
y \SI{0.80}{\second}, son menores que un periodo de muestreo y por tanto no distinguibles con los
datos disponibles.

La coincidencia entre ambos métodos depende de entregar los datos en variables de desviación: con
las señales absolutas la estimación converge a $K \approx \SI{1.47}{\celsius\per\percent}$ y
$\theta = 0$, valor que no es la ganancia del proceso sino el cociente
$T_1(\infty)/\Delta Q_1$, porque el estimador debe explicar también el nivel inicial de
temperatura.

\subsection{Validación cruzada y selección del modelo}

Cada modelo se simuló con la entrada del experimento que no se usó para identificarlo, manteniendo
fijos $K$, $\tau$ y $\theta$. La Tabla~\ref{tab:validacion} reporta las métricas calculadas desde
el primer cambio de $Q_1$, que es el intervalo que evalúa la respuesta dinámica, y la
Fig.~\ref{fig:validacion} muestra el cruce del modelo seleccionado junto con su residuo.

\begin{table}[!t]
\caption{Validación cruzada. Métricas calculadas desde el primer cambio de $Q_1$.}
\label{tab:validacion}
\centering
\footnotesize
\begin{tabular}{@{}lS[table-format=1.2]S[table-format=1.2]S[table-format=+1.2]S[table-format=1.3]S[table-format=2.1]@{}}
\toprule
\textbf{Modelo $\rightarrow$ datos} & {RMSE} & {MAE} & {Sesgo} & {$R^2$} & {Ajuste} \\
 & {(\si{\celsius})} & {(\si{\celsius})} & {(\si{\celsius})} & & {(\si{\percent})} \\
\midrule
Escalón $\rightarrow$ multinivel & 1.62 & 1.51 & +1.50 & 0.987 & 88.7 \\
Multinivel $\rightarrow$ escalón & 2.63 & 2.41 & -2.41 & 0.934 & 74.4 \\
\bottomrule
\end{tabular}
\end{table}

\begin{figure}[!t]
\centering
\begin{tikzpicture}
\begin{groupplot}[group style={group size=1 by 2, vertical sep=0.7cm}, estilo, height=3.6cm]
\nextgroupplot[ylabel={$T_1$ (\si{\celsius})}, xticklabels={}, legend pos=south east]
  \addplot[cAzul, aligera=5] table[col sep=comma, x=Tiempo_s, y=T1_experimental_C] {\G/validacionEscEnMul.csv};
  \addplot[cVerde, dashed, aligera=5] table[col sep=comma, x=Tiempo_s, y=T1_modelo_C] {\G/validacionEscEnMul.csv};
  \legend{Experimental, Modelo escalón}
\nextgroupplot[ylabel={Residuo (\si{\celsius})}, xlabel={Tiempo (\si{\second})}]
  \addplot[cRojo, aligera=5] table[col sep=comma, x=Tiempo_s, y=Residuo_C] {\G/validacionEscEnMul.csv};
  \addplot[cGris, dashed, domain=0:3520, samples=2] {0};
\end{groupplot}
\end{tikzpicture}
\caption{Validación del modelo identificado con el ensayo escalón sobre los datos multinivel, y
residuo de predicción $e(k) = T_{1,\mathrm{exp}}(k) - T_{1,\mathrm{mod}}(k)$.}
\label{fig:validacion}
\end{figure}

El modelo del ensayo escalón predice claramente mejor los datos independientes: su RMSE es un
\SI{38}{\percent} menor y su ajuste es \num{14.3} puntos porcentuales mayor. Ninguno de los dos
cruces está libre de sesgo. El modelo escalón subestima la temperatura en promedio
\SI{1.50}{\celsius}, y el residuo de la Fig.~\ref{fig:validacion} crece de forma ordenada con el
nivel de potencia: \num{+1.29}, \num{+0.88}, \num{+1.23}, \num{+1.85} y \SI{+2.20}{\celsius} para
$Q_1$ de \num{10} a \SI{50}{\percent}. El modelo multinivel presenta el sesgo opuesto,
\SI{-2.41}{\celsius}, es decir, sobreestima la temperatura del ensayo escalón. Ambos patrones
apuntan al mismo origen: la ganancia efectiva de la planta no es constante, de modo que un modelo
identificado en una región tiende a errar sistemáticamente en la otra. Esto es coherente con lo
observado en la Sección~II-A, donde el peso creciente del término radiativo hacía que una única
recta no describiera todo el rango.

Dado que ninguno de los dos modelos es válido en todo el rango, el criterio de selección fue la
región de operación prevista para el controlador. La referencia de trabajo es \SI{60}{\celsius},
que corresponde aproximadamente a $Q_1 = \SI{50}{\percent}$, justamente la condición del ensayo
escalón. Ese modelo además tuvo mejor capacidad predictiva sobre datos independientes. Su tiempo
muerto, $\theta = \SI{8.00}{\second}$, equivale a ocho periodos de muestreo, lo que permite
representar el retardo de forma exacta en la implementación discreta. El modelo nominal adoptado es
entonces

\begin{equation}
\Gpnom(s) = \frac{0{,}979\;e^{-8{,}00\,s}}{179{,}1\,s + 1}
\quad \left[\si{\celsius\per\percent}\right].
\label{eq:gpnom}
\end{equation}

\section{Diseño del controlador y evaluación MIL}

\subsection{Diseño del controlador PI}

Se empleó un controlador PI en forma paralela continua,

\begin{equation}
G_c(s) = \Kp + \frac{\Ki}{s} = \Kp\left(1 + \frac{1}{\Ti s}\right),
\qquad \Ki = \frac{\Kp}{\Ti},
\label{eq:pi}
\end{equation}

empleada en ambos diseños, con $e = \Tref - T_1$ en \si{\celsius} y salida en porcentaje de
potencia, de modo que
$[\Kp] = \si{\percent\per\celsius}$, $[\Ti] = \si{\second}$ y
$[\Ki] = \si{\percent\per\celsius\per\second}$. La acción integral es necesaria porque la planta no
tiene integrador propio y se requiere error estacionario nulo ante una referencia constante.

El primer diseño utilizó las reglas SIMC \cite{skogestad2003}, que para un modelo FOPDT dan

\begin{equation}
\Kp = \frac{1}{K}\,\frac{\tau}{\tau_c + \theta},
\qquad
\Ti = \min\{\tau,\; 4(\tau_c + \theta)\},
\label{eq:simc}
\end{equation}

donde $\tau_c$ es la constante de tiempo deseada en lazo cerrado. Se fijó $\tau_c = 2\theta$, un
valor conservador que prioriza robustez frente a rapidez. Con los parámetros de \eqref{eq:gpnom}
resulta $\tau_c = \SI{16}{\second}$ y $\tau_c + \theta = \SI{24}{\second}$; como
$4(\tau_c+\theta) = \SI{96}{\second} < \tau$, el tiempo integral queda truncado por la segunda rama
de \eqref{eq:simc}.

El segundo diseño se obtuvo con \emph{Control System Designer} sobre la misma planta nominal,
usando una estrategia distinta: en lugar de aplicar una fórmula de sintonización se moldeó la
respuesta en frecuencia del lazo abierto añadiendo un integrador y un cero real, y ajustando la
ganancia hasta cumplir dos criterios propios más exigentes que los de SIMC, margen de fase mayor a
\SI{70}{\degree} y margen de ganancia mayor a \SI{20}{\decibel}. La herramienta aproxima el retardo
mediante Padé para trazar el lugar de las raíces; los márgenes y las métricas reportadas se
recalcularon con el retardo exacto para evitar que esa aproximación contaminara la comparación.

Los dos controladores resultantes, en la forma \eqref{eq:pi}, son

\begin{equation}
G_{c,\mathrm{SIMC}}(s) = 7{,}62 + \frac{0{,}0794}{s},
\quad
G_{c,\mathrm{CSD}}(s) = 3{,}00 + \frac{0{,}0200}{s},
\label{eq:gc}
\end{equation}

con $\Kp$ en \si{\percent\per\celsius} y $\Ki$ en \si{\percent\per\celsius\per\second}. La
Tabla~\ref{tab:controladores} resume ambos diseños y su comportamiento sobre el modelo nominal.

\begin{table}[!t]
\caption{Controladores diseñados y su comportamiento sobre el modelo lineal nominal.}
\label{tab:controladores}
\centering
\footnotesize
\begin{tabular}{@{}lrr@{}}
\toprule
\textbf{Parámetro} & \textbf{PI--SIMC} & \textbf{PI--CSD} \\
\midrule
$\Kp$ (\si{\percent\per\celsius})                & 7,62   & 3,00 \\
$\Ki$ (\si{\percent\per\celsius\per\second})    & 0,0794 & 0,0200 \\
$\Ti$ (\si{\second})                             & 96     & 150 \\
\midrule
Frecuencia de cruce (\si{\radian\per\second})    & 0,0425 & 0,0168 \\
Margen de fase (\si{\degree})                     & 64,2   & 79,1 \\
Margen de ganancia (\si{\decibel})                & 13,3   & 21,5 \\
\midrule
$t_s$ lineal (\si{\second})                       & 202,3  & 161,9 \\
$M_p$ lineal (\si{\percent})                      & 7,76   & 1,67 \\
\bottomrule
\end{tabular}
\end{table}

Sobre el modelo lineal, sin restricciones de actuador, el controlador obtenido con
\emph{Control System Designer} se ve mejor en ambas métricas y además tiene márgenes mayores. Ese
resultado, sin embargo, no considera el límite de potencia, que es precisamente lo que evalúa la
etapa siguiente.

\subsection{Evaluación Model-in-the-Loop}

El controlador PI--SIMC de \eqref{eq:gc} se implementó de forma discreta con $T_s = \SI{1}{\second}$ e
integración por Euler hacia adelante,

\begin{equation}
I[k+1] = I[k] + \Ki T_s\, e[k],
\qquad
\Qcmd[k] = \Kp\, e[k] + I[k],
\label{eq:pidiscreto}
\end{equation}

distinguiendo la ganancia integral continua $\Ki$ del incremento por muestra $\Ki T_s$. La planta se
representó con \eqref{eq:gpnom} discretizada, con el retardo como ocho muestras de atraso. Se
simularon tres configuraciones bajo idénticas condiciones ---misma referencia, misma planta, misma
condición inicial, $\Tamb = \SI{24}{\celsius}$, $\Tref = \SI{60}{\celsius}$ y
\SI{1800}{\second} de simulación---: PI ideal sin restricciones, PI con saturación
$Q_1 \in [0,90]\,\si{\percent}$ sin anti-windup, y PI con saturación y anti-windup por
\emph{clamping}, que congela el integrador mientras la salida está saturada y el error empuja en
la dirección de la saturación \cite{mw_antiwindup}. Los resultados se muestran en la
Fig.~\ref{fig:mil} y la Tabla~\ref{tab:mil}; las métricas se calcularon con \texttt{stepinfo}
\cite{mw_stepinfo}.

\begin{figure*}[!t]
\centering
\begin{tikzpicture}
\begin{groupplot}[group style={group size=2 by 1, horizontal sep=1.3cm}, estilo, ancho,
                  xlabel={Tiempo (\si{\second})}, xmin=0, xmax=800]
\nextgroupplot[ylabel={$T_1$ (\si{\celsius})}, legend pos=south east, ymin=20, ymax=75]
  \addplot[cGris, dashed, domain=0:800, samples=2] {60};
  \addplot[cRojo]  table[col sep=comma, x index=0, y index=2] {\G/GRAFControladoresTemperatura.csv};
  \addplot[cVerde] table[col sep=comma, x index=0, y index=3] {\G/GRAFControladoresTemperatura.csv};
  \addplot[cAzul]  table[col sep=comma, x index=0, y index=1] {\G/GRAFControladoresTemperatura.csv};
  \legend{$\Tref$, Saturado sin AW, PI ideal, Saturado con AW}
\nextgroupplot[ylabel={Potencia (\si{\percent})}, legend pos=north east, ymin=-20, ymax=320]
  \addplot[cVerde]        table[col sep=comma, x index=0, y index=4] {\G/GRAFControladores.csv};
  \addplot[cRojo, dotted] table[col sep=comma, x index=0, y index=3] {\G/GRAFControladores.csv};
  \addplot[cRojo]         table[col sep=comma, x index=0, y index=2] {\G/GRAFControladores.csv};
  \addplot[cAzul]         table[col sep=comma, x index=0, y index=1] {\G/GRAFControladores.csv};
  \legend{$\Qcmd$ ideal, $\Qcmd$ sin AW, $Q_1$ sin AW, $Q_1$ con AW}
\end{groupplot}
\end{tikzpicture}
\caption{Evaluación MIL del PI--SIMC en tres configuraciones. Izquierda: temperatura. Derecha:
potencia comandada $\Qcmd$ y potencia aplicada $Q_1$. Se muestran los primeros
\SI{800}{\second} de los \SI{1800}{\second} simulados.}
\label{fig:mil}
\end{figure*}

\begin{table}[!t]
\caption{Métricas MIL. Las tres primeras columnas corresponden al PI--SIMC; la última, al PI--CSD
bajo idénticas condiciones de simulación.}
\label{tab:mil}
\centering
\footnotesize
\begin{tabular}{@{}lS[table-format=3.1]S[table-format=3.1]S[table-format=3.1]S[table-format=3.1]@{}}
\toprule
\textbf{Métrica} & {\textbf{Ideal}} & {\textbf{Sin AW}} & {\textbf{Con AW}} & {\textbf{CSD}} \\
\midrule
Tiempo de subida (\si{\second})        & 25.2  & 74.6  & 84.5  & 133.7 \\
$t_s$ (\si{\second})                   & 202.3 & 391.7 & 205.0 & 251.9 \\
$M_p$ (\si{\percent})                  & 7.76  & 30.1  & 0.00  & 0.00 \\
$e_{ss}$ (\si{\celsius})               & 0.00  & 0.00  & 0.00  & 0.00 \\
$T_{1,\max}$ (\si{\celsius})           & 62.8  & 70.8  & 60.0  & 60.0 \\
$Q_1$ mínimo aplicado (\si{\percent})  & 34.7  & 28.7  & 36.8  & 36.8 \\
$Q_1$ máximo aplicado (\si{\percent})  & 297.4 & 90.0  & 90.0  & 90.0  \\
Tiempo con $Q_1$ saturada (\si{\second}) & {---} & 127   & 66    & 21 \\
\bottomrule
\end{tabular}
\end{table}

El PI ideal cumple los requisitos en simulación, pero exige un máximo de \SI{297}{\percent} de
potencia. Ese resultado es físicamente irrealizable y muestra que evaluar el controlador sin las
restricciones del actuador no permite concluir nada sobre su comportamiento real.

Al introducir la saturación sin anti-windup el desempeño se degrada de forma marcada: el
sobreimpulso pasa de \SI{7.76}{\percent} a \SI{30.1}{\percent} y el tiempo de establecimiento casi
se duplica. El mecanismo es visible en la Fig.~\ref{fig:mil}. Durante los primeros
\SI{127}{\second} el comando permanece por encima del límite, de modo que el error sigue
integrándose aunque la potencia aplicada ya no pueda aumentar. Cuando $T_1$ cruza la referencia en
$t = \SI{103}{\second}$ el integrador acumula un excedente que el error negativo debe descargar
antes de que el comando vuelva al rango útil, lo que ocurre \SI{24}{\second} después. La planta
entretanto sigue recibiendo el \SI{90}{\percent} de potencia y alcanza \SI{70.8}{\celsius}. Para
recuperarse, el controlador debe reducir la potencia hasta el \SI{28.7}{\percent} en
$t = \SI{186}{\second}$, por debajo del \SI{36.8}{\percent} que la planta necesita en régimen.

Con \emph{clamping} el integrador deja de acumular mientras la salida está saturada, la saturación
dura \SI{66}{\second} en lugar de \SI{127}{\second} y la respuesta llega a la referencia sin
sobrepasarla. De las tres configuraciones, solo esta cumple simultáneamente
$t_s \leq \SI{600}{\second}$ y $M_p \leq \SI{5}{\percent}$; la configuración sin anti-windup falla
el requisito de sobreimpulso por un factor de seis. Conviene notar que el anti-windup cuesta
\SI{10}{\second} de tiempo de subida frente al caso sin protección, un precio bajo comparado con lo
que evita.

La última columna de la Tabla~\ref{tab:mil} incorpora el PI--CSD con saturación y anti-windup, bajo
la misma planta, el mismo $T_s$, los mismos límites, la misma condición inicial y el mismo tiempo de
simulación.\footnote{En esa corrida la columna PI--SIMC da $t_r = \SI{84.56}{\second}$ y
$t_s = \SI{205.18}{\second}$; la diferencia respecto de la tercera columna es de interpolación
sub-muestra de \texttt{stepinfo} y no de comportamiento.} Ambos controladores cumplen los
requisitos y ambos alcanzan la referencia sin sobreimpulso y sin error estacionario, de modo que la
comparación se decide por rapidez y por esfuerzo de control. El PI--SIMC establece \SI{47}{\second}
antes, pero pide un comando inicial de \SI{274}{\percent} y mantiene el actuador saturado
\SI{66}{\second}; el PI--CSD pide \SI{108}{\percent} y satura solo \SI{21}{\second}, a cambio de un
tiempo de subida \SI{57}{\percent} mayor. Es decir, el PI--CSD opera mucho más cerca de la
capacidad real del actuador y tiene márgenes de estabilidad mayores, mientras que el PI--SIMC
aprovecha mejor la autoridad disponible.

Ninguno de los dos domina al otro. Como ambos satisfacen las especificaciones con holgura y el
tiempo de establecimiento no es crítico en esta planta, un criterio de robustez favorecería al
PI--CSD. Para las etapas de generación y verificación de código se utilizó el PI--SIMC con
saturación y anti-windup, según lo establecido en la guía de laboratorio, y el PI--CSD se conservó
como contraste metodológico.

\section{Generación de código y verificación de software}

\subsection{Implementación del controlador}

El sistema bajo prueba (\emph{System Under Test}, SUT) es el modelo \texttt{Controller.slx}, que
contiene únicamente el algoritmo de control. Su interfaz es
$\Tref,\,T_1 \rightarrow Q_1$: recibe la referencia y la temperatura medida, ambas en
\si{\celsius}, y entrega la potencia comandada en porcentaje. Internamente ejecuta la ley discreta
\eqref{eq:pidiscreto} con $T_s = \SI{1}{\second}$, aplica la saturación
$Q_1 \in [0,90]\,\si{\percent}$ y el anti-windup por \emph{clamping}. La Fig.~\ref{fig:controller}
muestra su estructura.

\pendiente{Insertar aquí la captura de \texttt{Controller.slx} (Simulink) como
\texttt{figuras/controller.png}. Es la única captura que el instructivo pide explícitamente.}

\begin{figure}[!t]
\centering
% \includegraphics[width=\columnwidth]{figuras/controller.png}
\framebox[0.95\columnwidth][c]{\rule{0pt}{3.2cm}\scriptsize Figura de \texttt{Controller.slx}}
\caption{Modelo \texttt{Controller.slx}: implementación discreta del PI--SIMC con saturación y
anti-windup por \emph{clamping}.}
\label{fig:controller}
\end{figure}

La planta no forma parte de este modelo. Esa separación es deliberada: es lo que permite generar
código solo del algoritmo, instanciarlo varias veces dentro de un banco de pruebas y sustituir una
instancia por su versión compilada sin alterar el resto del lazo. Para ello el modelo se configuró
con \texttt{Code interface packaging = Reusable function} y se habilitaron múltiples instancias por
modelo de nivel superior.

De la configuración de generación de código solo interesan aquí tres decisiones. La primera es el
uso de \texttt{ert.tlc} con lenguaje C bajo el estándar C99, que produce código orientado a sistemas
embebidos \cite{mw_codegen}. La segunda es el solver de paso fijo discreto con paso $T_s$, coherente
con \eqref{eq:pidiscreto}. La tercera es \texttt{Default parameter behavior = Inlined} combinado con
la clase de almacenamiento \texttt{ExportedGlobal} para $Kp_{SIMC}$ y $Ki_{SIMC}$: el resto de las
constantes se incrusta en el código, pero las dos ganancias permanecen como variables globales y por
tanto siguen siendo calibrables sin regenerar. Para la etapa SIL el hardware objetivo fue la propia
computadora (Intel x86-64), no todavía el Arduino.

Las comprobaciones previas a la generación ---\emph{Code Generation Advisor} con los objetivos
\emph{Execution efficiency} y \emph{Traceability}, y \emph{Model Advisor}--- no reportaron ningún
fallo. La única advertencia relevante se refiere al soporte de valores no finitos en el entorno de
software, que no afecta al algoritmo implementado; no se aplicaron correcciones automáticas para
eliminarla.

\subsection{Inspección del código generado}

La revisión del código C verificó que las características esenciales del controlador sobrevivieran
a la generación. La Tabla~\ref{tab:codigo} resume la correspondencia encontrada.

\begin{table}[!t]
\caption{Correspondencia entre el modelo y el código generado.}
\label{tab:codigo}
\centering
\footnotesize
\begin{tabular}{@{}p{0.27\columnwidth}p{0.63\columnwidth}@{}}
\toprule
\textbf{Característica} & \textbf{Evidencia en el código} \\
\midrule
Interfaz &
\texttt{void Controller(const real\_T *rtu\_T\_ref, const real\_T *rtu\_T1, real\_T *rty\_Q1,
DW\_Controller\_f\_T *localDW)} \\
\addlinespace
Estado del integrador &
Campo único \texttt{Integrator\_DSTATE} dentro de la estructura de estado \texttt{DW\_Controller\_f\_T} \\
\addlinespace
Parámetros calibrables &
\texttt{real\_T Kp\_SIMC = 7.620337;} y \texttt{real\_T Ki\_SIMC = 0.079379;} definidos como
variables globales \\
\addlinespace
Saturación &
Ramas explícitas con los límites \texttt{90.0} y \texttt{0.0} sobre la salida \\
\addlinespace
Anti-windup &
Condición que anula el incremento integral cuando hay exceso y su signo coincide con el del error \\
\addlinespace
Tipo de dato &
\texttt{typedef double real\_T} \\
\addlinespace
Separación controlador--planta &
El archivo generado no contiene ningún modelo de la planta \\
\bottomrule
\end{tabular}
\end{table}

El punto menos evidente es el anti-windup, porque en el código no aparece como un bloque sino como
una condición aritmética. El fragmento correspondiente es

\begin{lstlisting}
if ((rtb_DeadZone != 0.0) && (tmp == tmp_0)) {
  rtb_IntegralGain = 0.0;
}
localDW->Integrator_DSTATE += rtb_IntegralGain;
\end{lstlisting}

La variable \texttt{rtb\_DeadZone} mide cuánto excede el comando el rango permitido y vale cero
mientras no hay saturación; \texttt{tmp} y \texttt{tmp\_0} codifican los signos del exceso y del
término integral. Cuando hay saturación y ambos signos coinciden, el incremento se fuerza a cero y
el estado del integrador no cambia. Esa es exactamente la lógica de \emph{clamping} evaluada en
MIL.

\subsection{Verificación Software-in-the-Loop}

En SIL el código generado se compila y ejecuta sobre la computadora dentro de la misma simulación
\cite{mw_sil}. El banco \texttt{TCLab\_Test\_SIL.slx} contiene dos ramas que comparten referencia,
temperatura ambiente, modelo de planta, condición inicial, periodo de muestreo y tiempo de
simulación; la única diferencia es que una instancia del controlador se ejecuta en modo Normal y la
otra en modo SIL. La comparación es por tanto \emph{back-to-back}: misma entrada, dos
implementaciones del mismo algoritmo.

Las diferencias evaluadas fueron

\begin{equation}
\Delta Q_1 = Q_{1,\mathrm{SIL}} - Q_{1,\mathrm{Normal}},
\qquad
\Delta T_1 = T_{1,\mathrm{SIL}} - T_{1,\mathrm{Normal}},
\label{eq:difsil}
\end{equation}

sobre las \num{1801} muestras de los \SI{1800}{\second} simulados. Ambas resultaron idénticamente
nulas, como se observa en la Fig.~\ref{fig:diferencias} y se cuantifica en la
Tabla~\ref{tab:silpil}, de modo que la prueba satisface el criterio de equivalencia con igualdad
exacta y no solo dentro de una tolerancia.

Este resultado verifica que la traducción de modelo a código C no introdujo diferencias numéricas:
el orden de las operaciones, la actualización del estado y la lógica de saturación y anti-windup se
conservaron. No verifica, en cambio, nada relacionado con el procesador final, porque el código se
compiló para la arquitectura del computador anfitrión y con aritmética \texttt{double}. Tampoco
dice nada sobre la planta, que sigue siendo un modelo.

\section{Verificación Processor-in-the-Loop}

En PIL el código del controlador se compila para el procesador objetivo, se descarga a la tarjeta y
se ejecuta allí, mientras Simulink intercambia con ella las señales de entrada y salida en cada
periodo de muestreo \cite{mw_pil}. El procesador empleado fue un Arduino Uno R4 Minima, utilizado
exclusivamente como plataforma de verificación y no como el destino final del controlador.

El reparto entre lo que se ejecuta físicamente y lo que permanece simulado es el punto clave de esta
etapa. Sobre el Arduino se ejecuta únicamente el algoritmo de control: la ley PI discreta, la
saturación y el anti-windup, con el estado del integrador residiendo en la memoria del
microcontrolador. Todo lo demás sigue siendo un modelo dentro de Simulink: la planta FOPDT, el
tiempo muerto, la temperatura ambiente y la referencia. No intervienen el sensor ni el calentador de
la TCLab.

El banco \texttt{TCLab\_Test\_PIL.slx} emplea de nuevo dos ramas que referencian
\texttt{Controller.slx}, una en modo Normal y otra en modo PIL. A diferencia del banco SIL, las dos
ramas comparten solo las fuentes $\Tref$ y $\Tamb$, y cada una conserva su propia copia del retardo
y de la planta. Esa separación es necesaria porque, si las ramas compartieran la planta, cualquier
discrepancia entre las acciones de control produciría trayectorias de temperatura distintas y
dejaría de poder atribuirse la diferencia a una sola causa. Tampoco se mantuvieron simultáneamente
una instancia SIL y una PIL referenciando el mismo modelo.

Las diferencias evaluadas fueron las análogas a \eqref{eq:difsil},
$\Delta Q_1 = Q_{1,\mathrm{PIL}} - Q_{1,\mathrm{Normal}}$ y
$\Delta T_1 = T_{1,\mathrm{PIL}} - T_{1,\mathrm{Normal}}$, sobre las mismas \num{1801} muestras. El
resultado fue también de igualdad exacta (Fig.~\ref{fig:diferencias} y Tabla~\ref{tab:silpil}). No
se observó ninguna diferencia que analizar, lo cual es consistente con haber conservado el tipo
\texttt{double} y la misma estructura algorítmica en las tres ejecuciones.

\begin{figure}[!t]
\centering
\begin{tikzpicture}
\begin{groupplot}[group style={group size=1 by 2, vertical sep=0.7cm}, estilo, height=3.3cm,
                  xmin=0, xmax=1800, legend pos=north east]
\nextgroupplot[ylabel={$\Delta Q_1$ (\si{\percent})}, xticklabels={}, ymin=-1, ymax=1]
  \addplot[cAzul, line width=1.4pt, opacity=0.45] table[col sep=comma, x index=0,
    y expr=\thisrowno{2}-\thisrowno{1}] {\G/GRAFComparacionNormalSIL.csv};
  \addplot[cNaranja, dashed] table[col sep=comma, x index=0,
    y expr=\thisrowno{2}-\thisrowno{1}] {\G/GRAFComparacionNormalPIL.csv};
  \legend{SIL, PIL}
\nextgroupplot[ylabel={$\Delta T_1$ (\si{\celsius})}, xlabel={Tiempo (\si{\second})},
               ymin=-1, ymax=1]
  \addplot[cAzul, line width=1.4pt, opacity=0.45] table[col sep=comma, x index=0,
    y expr=\thisrowno{4}-\thisrowno{3}] {\G/GRAFComparacionNormalSIL.csv};
  \addplot[cNaranja, dashed] table[col sep=comma, x index=0,
    y expr=\thisrowno{4}-\thisrowno{3}] {\G/GRAFComparacionNormalPIL.csv};
  \legend{SIL, PIL}
\end{groupplot}
\end{tikzpicture}
\caption{Diferencias \emph{back-to-back} respecto de la ejecución Normal. Las cuatro trazas son
idénticamente nulas en las \num{1801} muestras; la escala vertical se fijó en
$\pm\num{1}$ para hacer visible que no hay desviación alguna.}
\label{fig:diferencias}
\end{figure}

\begin{table}[!t]
\caption{Resultados cuantitativos de las verificaciones SIL y PIL, sobre \num{1801} muestras
($T_s = \SI{1}{\second}$, \SI{1800}{\second} de simulación).}
\label{tab:silpil}
\centering
\footnotesize
\begin{tabular}{@{}lcccccc@{}}
\toprule
& \multicolumn{2}{c}{\textbf{SIL}} & \multicolumn{2}{c}{\textbf{PIL}} & \multicolumn{2}{c}{\textbf{Igualdad exacta}} \\
\cmidrule(lr){2-3}\cmidrule(lr){4-5}\cmidrule(lr){6-7}
\textbf{Señal} & Máx.\ abs. & RMSE & Máx.\ abs. & RMSE & SIL & PIL \\
\midrule
$Q_1$ (\si{\percent})   & 0 & 0 & 0 & 0 & Sí & Sí \\
$T_1$ (\si{\celsius})   & 0 & 0 & 0 & 0 & Sí & Sí \\
\midrule
\multicolumn{5}{@{}l}{\textbf{Resultado}} & \textbf{PASS} & \textbf{PASS} \\
\bottomrule
\end{tabular}
\end{table}

La evidencia que aporta PIL es distinta de la de SIL. SIL comprueba la traducción de modelo a código
usando el compilador y la aritmética del computador anfitrión; PIL añade que ese mismo código se
compila con la cadena de herramientas del procesador objetivo, cabe en su memoria y produce los
mismos resultados ejecutándose sobre él. Es decir, PIL cubre los efectos del compilador cruzado, de
los tamaños de palabra y del comportamiento aritmético del dispositivo, que SIL no puede detectar.

Aun así, un PASS en PIL no constituye una validación del controlador sobre la TCLab. Durante la
prueba la planta y el retardo siguen siendo modelos, $\Tamb$ es una señal de Simulink, no se mide
temperatura con el sensor real ni se acciona el calentador, no se ejercitan los canales físicos de
entrada y salida, y la ejecución no se realiza en tiempo real contra un proceso físico. Lo que PIL
verifica es el código y su ejecución sobre el procesador, no el lazo cerrado sobre la planta.

\section{Implementación experimental}

\subsection{Arquitectura del lazo físico}

En esta etapa desaparece el modelo de la planta. El lazo implementado es

\begin{center}
$\Tref \rightarrow$ controlador (Arduino Leonardo) $\rightarrow Q_1 \rightarrow$ calentador
$\rightarrow T_1 \rightarrow$ controlador,
\end{center}

y se ejecutó sobre el Arduino Leonardo integrado en la TCLab mediante Monitor \& Tune
\cite{mw_arduino}. El algoritmo del controlador se copió a
\texttt{Controller\_implementacion.slx} para no alterar \texttt{Controller.slx}, que se conservó
como referencia de las etapas de verificación. La única diferencia entre ambos es la representación
numérica: en la versión de implementación las entradas, las ganancias, los límites y el estado del
integrador se resolvieron en \texttt{single}, adecuado para el microcontrolador de 8 bits, mientras
que la ley de control, la saturación, el anti-windup y $T_s = \SI{1}{\second}$ se mantuvieron sin
cambios.

La cadena de señal es la siguiente. La temperatura se lee del canal analógico \texttt{A0}, se
convierte a \texttt{single} y se escala mediante una ganancia $500/1023$ seguida de un
desplazamiento de \SI{-50}{\celsius}, correspondientes al rango del sensor y a la resolución del
convertidor de 10 bits. La medición pasa por un filtro de media móvil con $\lambda = 0{,}5$ antes de
entrar al controlador. La salida se satura, se escala a la resolución del PWM y se convierte a
\texttt{uint8} para el pin \texttt{D3}, que comanda el calentador.

Es importante precisar el papel de Monitor \& Tune. Al iniciarlo, Simulink genera el código del
modelo y sus modelos referenciados, lo compila y lo descarga al Arduino
\cite{mw_arduino_workflows}. Todas las operaciones del lazo ---adquisición de $T_1$, conversión,
filtrado, ejecución del PI, cálculo y limitación de $Q_1$ y generación del PWM--- se ejecutan sobre
el Arduino Leonardo. La computadora no participa en el cálculo del lazo de control: el enlace de
\emph{External Mode} sirve para monitorizar señales, registrar resultados y ajustar parámetros
durante la ejecución.

\subsection{Resultados experimentales}

Tras una ejecución previa con $\Tref = \SI{25}{\celsius}$ para comprobar la lectura del sensor y
los límites del actuador, el experimento definitivo se realizó con $\Tref = \SI{60}{\celsius}$
durante \SI{607}{\second} (\num{608} muestras), partiendo de $T_1 = \SI{26.25}{\celsius}$. La Fig.~\ref{fig:experimental} presenta las señales registradas junto
con la respuesta MIL, y la Tabla~\ref{tab:experimental} compara ambas cuantitativamente.

\begin{figure}[!t]
\centering
\begin{tikzpicture}
\begin{groupplot}[group style={group size=1 by 2, vertical sep=0.7cm}, estilo, height=3.6cm,
                  xmin=0, xmax=607]
\nextgroupplot[ylabel={Temperatura (\si{\celsius})}, xticklabels={}, ymin=22, ymax=66,
               legend pos=south east]
  \addplot[cGris, dashed] table[col sep=comma, x index=0, y index=1] {\G/ImplementacionA12.csv};
  \addplot[cAzul]         table[col sep=comma, x index=0, y index=2] {\G/ImplementacionA12.csv};
  \addplot[cVerde, dotted, line width=1pt]
    table[col sep=comma, x index=0, y index=1] {\G/GRAFControladoresTemperatura.csv};
  \legend{$\Tref$, $T_1$ experimental, $T_1$ MIL}
\nextgroupplot[ylabel={$Q_1$ (\si{\percent})}, xlabel={Tiempo (\si{\second})}, ymin=0, ymax=100,
               legend pos=north east]
  \addplot[cNaranja] table[col sep=comma, x index=0, y index=3] {\G/ImplementacionA12.csv};
  \addplot[cVerde, dotted, line width=1pt]
    table[col sep=comma, x index=0, y index=1] {\G/GRAFControladores.csv};
  \legend{Experimental, MIL}
\end{groupplot}
\end{tikzpicture}
\caption{Implementación sobre la TCLab con $\Tref = \SI{60}{\celsius}$, superpuesta a la respuesta
MIL del mismo controlador.}
\label{fig:experimental}
\end{figure}

\begin{table}[!t]
\caption{Desempeño experimental frente al obtenido en MIL. Las etapas SIL y PIL reprodujeron
exactamente el resultado MIL.}
\label{tab:experimental}
\centering
\footnotesize
\begin{tabular}{@{}lS[table-format=3.1]S[table-format=3.2]@{}}
\toprule
\textbf{Métrica} & {\textbf{MIL}} & {\textbf{Experimental}} \\
\midrule
Tiempo de subida (\si{\second})           & 84.5  & 74 \\
$t_s$ (\si{\second})                      & 205.0 & 298 \\
$M_p$ (\si{\percent})                     & 0.00  & 1.96 \\
$e_{ss}$ (\si{\celsius})                  & 0.00  & 0.04 \\
$T_{1,\max}$ (\si{\celsius})              & 60.0  & 60.66 \\
$Q_1$ mínimo aplicado (\si{\percent})     & 36.8  & 25.42 \\
$Q_1$ en régimen (\si{\percent})          & 36.8  & 40.3 \\
Tiempo con $Q_1$ saturada (\si{\second})  & 66    & 69 \\
\bottomrule
\end{tabular}
\end{table}

El controlador cumple su función sobre la planta real: alcanza la referencia, la mantiene y respeta
en todo momento el intervalo de potencia permitido, con un máximo del \SI{90}{\percent} durante el
calentamiento inicial y sin saturación sostenida injustificada. El error estacionario, calculado como la diferencia entre la
referencia y la temperatura media para $t \geq \SI{300}{\second}$, es de \SI{0.04}{\celsius}.
Ese valor equivale a menos de una décima parte del escalón mínimo del convertidor,
$500/1023 = \SI{0.489}{\celsius}$; es decir, el error queda muy por debajo de la resolución con que
el sistema puede medir la temperatura.

Las diferencias respecto de MIL son sistemáticas y admiten interpretación. El sobreimpulso de
\SI{1.96}{\percent} merece una precisión: el máximo absoluto del registro,
\SI{60.66}{\celsius}, no ocurre durante el transitorio sino en $t = \SI{520}{\second}$, ya en
régimen permanente, mientras que el máximo alcanzado antes de que la respuesta se
establezca, en $t = \SI{273}{\second}$, es de \SI{59.80}{\celsius}, por debajo de la referencia. Lo que la métrica está recogiendo es la
oscilación de baja amplitud alrededor del punto de operación ---la desviación típica de $T_1$ para
$t \geq \SI{300}{\second}$ es de \SI{0.245}{\celsius}--- y no un sobredisparo de la respuesta al
escalón.

La discrepancia más informativa está en el régimen permanente. En MIL el lazo se estabiliza con
$Q_1 = \SI{36.8}{\percent}$, mientras que la planta real necesita \SI{40.3}{\percent} para sostener
la misma temperatura. La ganancia estática aparente durante el experimento resulta entonces de
aproximadamente \SI{0.84}{\celsius\per\percent}, cerca de un \SI{15}{\percent} por debajo del
$K = \SI{0.979}{\celsius\per\percent}$ del modelo nominal. Este error de ganancia explica también el
tiempo de establecimiento mayor: con menos ganancia efectiva de planta, el lazo cierra más lento,
y el sistema tarda \SI{298}{\second} en lugar de \SI{205}{\second} en entrar en la banda del
\SI{2}{\percent}. El tiempo de subida, en cambio, resultó menor que en MIL (\SI{74}{\second} frente
a \SI{84.5}{\second}), lo que es coherente con haber partido de \SI{26.25}{\celsius} en lugar de
\SI{24}{\celsius}: el escalón real fue un \SI{6}{\percent} más pequeño.

Otras diferencias visibles en la Fig.~\ref{fig:experimental} no aparecen en MIL por construcción.
La señal $Q_1$ experimental es ruidosa en régimen (desviación típica de \SI{1.81}{\percent}) porque
el controlador amplifica el ruido de medición por $\Kp = \SI{7.62}{\percent\per\celsius}$, de modo
que un solo escalón del convertidor produce una variación de casi cuatro puntos porcentuales de
potencia. El descenso hasta el \SI{25.4}{\percent} en $t = \SI{115}{\second}$, más
profundo que el \SI{36.8}{\percent} de MIL, corresponde a la corrección tras la saturación inicial
sobre una planta cuya dinámica real no coincide exactamente con el FOPDT identificado.

\section{Discusión}

El recorrido completo deja ver que cada etapa del flujo responde una pregunta distinta, y que
ninguna sustituye a las demás.

El modelo obtenido del balance de energía resultó útil para entender el sistema pero no para
diseñar. Su linealización es válida solo en un entorno reducido del punto de operación, como
cuantifica la Tabla~\ref{tab:fisico}, mientras que el transitorio que debe manejar el controlador
es precisamente un calentamiento desde temperatura ambiente. Por eso el modelo de diseño tuvo que
obtenerse de forma experimental. El FOPDT identificado con el ensayo escalón cumple
esa función, pero la validación cruzada mostró también su límite: los residuos no son ruido
alrededor de cero sino un patrón que crece con el nivel de potencia. Esto significa que el modelo no
falla al azar, sino que subestima la temperatura de forma creciente a medida que se aleja de la
región donde fue identificado, y que un solo modelo lineal con retardo no describe la planta en todo
el rango. La elección del modelo nominal fue, en ese sentido, una elección de región de operación
más que de calidad absoluta.

El efecto de la saturación resultó ser el factor de diseño más relevante de toda la práctica. El PI
sintonizado con SIMC demanda un \SI{274}{\percent} de potencia en el primer instante, casi tres
veces lo que el actuador puede entregar. Con el actuador limitado pero sin protección del
integrador, el controlador sintonizado correctamente incumple el requisito de sobreimpulso por un
factor de seis. La causa no está en la sintonización sino en que el integrador sigue acumulando
mientras la potencia ya no puede crecer, de modo que el estado del controlador deja de corresponder
a lo que realmente ocurre en la planta. El \emph{clamping} corrige eso deteniendo la acumulación
durante la saturación, y el costo es de apenas \SI{10}{\second} de tiempo de subida. La lección es
que el desempeño de un controlador no puede evaluarse sin las restricciones del actuador: el caso
ideal de la Tabla~\ref{tab:mil} cumple los requisitos y a la vez es irrealizable.

La comparación entre los dos diseños refuerza esa idea. Sobre el modelo lineal sin restricciones el
PI--CSD parece claramente superior, con mejor tiempo de establecimiento, menor sobreimpulso y
márgenes mayores. Al incorporar el límite de potencia el orden cambia parcialmente: el PI--SIMC
resulta más rápido porque aprovecha toda la autoridad disponible, mientras que el PI--CSD se
mantiene mucho más cerca de la capacidad real del actuador. Ambos cumplen, y decidir entre ellos
exige un criterio explícito ---rapidez frente a robustez y esfuerzo de control--- que no se deduce
de las métricas por sí solas.

Las verificaciones SIL y PIL cumplieron un papel distinto del de MIL. MIL evalúa si el algoritmo
resuelve el problema de control; SIL y PIL no evalúan desempeño en absoluto, sino conservación de
comportamiento al cambiar de representación. Que ambas dieran igualdad exacta y no solo equivalencia
dentro de una tolerancia da confianza en que las diferencias observadas después, sobre la planta
real, no provienen del código ni del procesador. Ese es su valor práctico: al llegar a la
implementación física, el conjunto de causas posibles de una discrepancia ya está acotado al modelo
de la planta y al hardware de entrada y salida. Al mismo tiempo, conviene no sobreinterpretar el
resultado: durante PIL la planta, el retardo y la temperatura ambiente siguen siendo señales de
Simulink, y no se ejercita ningún canal físico.

Efectivamente, las diferencias aparecieron al cerrar el lazo sobre la TCLab, y son del tipo que las
etapas anteriores no podían revelar. La más significativa es el error de ganancia: la planta
necesita en régimen más potencia que la prevista (Tabla~\ref{tab:experimental}), lo que equivale a
una ganancia estática aparente cerca de un \SI{15}{\percent} menor que la nominal. Esa
diferencia es coherente con el sesgo detectado en la validación cruzada y con la dependencia de la
ganancia respecto de la región de operación, y se traduce en un lazo más lento que el simulado. A
ello se suman efectos ausentes del modelo: el ruido de medición amplificado por
$\Kp = \SI{7.62}{\percent\per\celsius}$, la cuantización del convertidor, y variaciones ambientales
durante los \SI{607}{\second} del experimento. Ninguno de estos efectos invalida el diseño ---el
controlador cumple los requisitos sobre la planta real--- pero sí muestran que el margen con que se
cumplen es menor que el que sugería la simulación, y que ese margen proviene en buena medida de
haber elegido un $\tau_c$ conservador.

\section{Conclusiones}

\begin{enumerate}[leftmargin=*]

\item El modelo FOPDT identificado con el ensayo escalón,
$\Gpnom(s) = 0{,}979\,e^{-8{,}00 s}/(179{,}1 s + 1)$, representa adecuadamente la planta en la
región de operación del controlador, con un ajuste del \SI{88.7}{\percent} sobre datos no empleados
en su identificación, frente al \SI{74.4}{\percent} del modelo multinivel. Su limitación es
igualmente clara: el residuo de validación no es aleatorio sino creciente con el nivel de potencia,
de \SI{+0.88}{\celsius} a \SI{+2.20}{\celsius} entre $Q_1 = \SI{20}{\percent}$ y
\SI{50}{\percent}, lo que indica que una sola ganancia no describe todo el rango de la planta.

\item El anti-windup, y no la sintonización, fue lo que determinó el cumplimiento de las
especificaciones. Con el mismo $\Kp$ y el mismo $\Ti$ obtenidos por SIMC, la configuración sin
protección del integrador produce $M_p = \SI{30.1}{\percent}$ e incumple el requisito
$M_p \leq \SI{5}{\percent}$ por un factor de seis, mientras que con \emph{clamping} se obtiene
$M_p = \SI{0}{\percent}$ y $t_s = \SI{205}{\second}$. Haber aplicado correctamente las reglas SIMC
no demuestra por sí mismo que el controlador cumpla los requisitos.

\item Los dos diseños PI cumplen las especificaciones en MIL, pero con compromisos opuestos y sin
que ninguno domine al otro: el PI--SIMC establece \SI{47}{\second} antes, y el PI--CSD demanda un
comando inicial de \SI{108}{\percent} frente a \SI{274}{\percent} y mantiene el actuador saturado
\SI{21}{\second} frente a \SI{66}{\second}, además de contar con márgenes de estabilidad mayores.
La selección depende del criterio que se privilegie y no se deduce de las métricas de desempeño por
sí solas.

\item Las verificaciones SIL y PIL resultaron en igualdad exacta de $Q_1$ y $T_1$ sobre
\num{1801} muestras, lo que acredita que ni la generación de código ni su ejecución sobre el
Arduino Uno R4 Minima alteraron el comportamiento del controlador validado en MIL. Ese resultado no
constituye una validación del sistema completo: durante ambas pruebas la planta, el retardo y la
temperatura ambiente permanecen simulados y no se ejercitan los canales físicos de la TCLab, de
modo que el valor de estas etapas está en acotar las causas posibles de una discrepancia
posterior, no en anticiparla.

\item La implementación sobre la TCLab cumple los requisitos ---$t_s = \SI{298}{\second}$ y
$M_p = \SI{1.96}{\percent}$--- con un error estacionario de \SI{0.04}{\celsius}, muy inferior al
escalón mínimo del convertidor (\SI{0.489}{\celsius}). Las diferencias respecto de MIL se explican
principalmente por un error de ganancia del modelo: la planta requiere \SI{40.3}{\percent} de
potencia en régimen frente al \SI{36.8}{\percent} previsto, lo que corresponde a una ganancia
estática aparente cerca de un \SI{15}{\percent} menor que la nominal y explica el tiempo de
establecimiento un \SI{45}{\percent} mayor.

\end{enumerate}


\balance
\bibliographystyle{IEEEtran}
\bibliography{referencias}

\end{document}
